
\subsubsection{Interpretation}

The extraction bottleneck indicates either implementation inefficiency correctable through optimization or fundamental computational limits inherent to 7B model scale. Modern inference libraries (FlashAttention-2, vLLM) achieve 10-100× speedups over naive implementations through kernel fusion and memory optimization. If the overhead stems from unoptimized forward passes and tensor operations, optimized libraries could reduce per-example cost from 12 minutes to under 1 minute.

Alternatively, the bottleneck may reflect inherent computational demands. Hidden state extraction requires 246 forward passes through a 7B parameter model with 32 layers and attention mechanisms. Each forward pass involves billions of parameter operations. The 15.4 MB tensor storage requirement is small compared to the 2.4GB model memory, suggesting computation rather than memory dominates the bottleneck.

The root cause—implementation inefficiency versus fundamental scaling—remains unresolved without profiling studies measuring time breakdown across forward pass, tensor extraction, and I/O operations.

\subsubsection{Limitations}

\#\#\#\# Zero Empirical Evidence

No correlation measurements were obtained. All predictions (P1, P2, P3) remain untested. The hypothesis validity is unknown. However, documenting methodological barriers provides value by informing resource allocation decisions for future geometric uncertainty research.

\#\#\#\# Model Substitution

Mistral-7B-v0.1 was substituted for the planned Llama-3-8B-Instruct due to access restrictions. Both are 7B-scale decoder-only transformers with 32 layers and hidden dimension 4096. However, Mistral uses sliding window attention while Llama-3 uses full attention, potentially affecting hidden state geometry. The computational bottleneck findings apply regardless of specific model choice, as both require similar tensor operations.

\#\#\#\# Arbitrary Layer Range

Layers 24-31 were selected based on proximity to output logits (hypothesis: late layers encode decision uncertainty) without empirical validation. No ablation study tested whether earlier layers (15-22), later layers (28-31), or single-layer extraction (layer 31) would achieve better correlation or lower computational cost. Layer selection represents an untested assumption.

\#\#\#\# Efficiency Claim Contradiction

The original hypothesis claimed geometric features would enable sub-10ms production overhead compared to semantic entropy's approximately 1500ms. The extraction bottleneck contradicts this: if extraction takes 12 minutes per example (observed) versus semantic entropy's approximately 1.5 seconds per example (from literature), the geometric approach is 480-fold slower, not faster. The original efficiency claim assumed extraction overhead is negligible (single forward pass ≈ 100ms). Discovering that extraction is the bottleneck revises understanding of computational requirements.

\subsubsection{Computational Planning}

The 20-fold complexity underestimate reveals a limitation in hypothesis planning methodology. Phase 3 complexity scores rely on intuitive estimation without empirical validation. For tensor-heavy operations on large models, intuition systematically underestimates resource requirements.

Future Phase 3 planning should include empirical profiling on small subsets (N=10) before assigning complexity scores to operations involving large tensor extractions, iterative optimization on 7B+ models, or multi-sample generation. Profiling overhead (approximately 30 minutes) is small compared to cost of 10-hour failed experiments.
